% Source: https://artofproblemsolving.com/community/c3h1985059_2018_amc_10b_problem_6
% Stable identity: aops-topic-1985059-q01
Imagine that the chips are drawn in a uniformly random order and look at the first two. All $5\cdot4=20$ ordered pairs are equally likely. Three draws are needed exactly when the first two chips have sum at most $4$, i.e. they are $(1,2),(2,1),(1,3),(3,1)$. So the probability is
\[
\dfrac{4}{20}=\boxed{\dfrac{1}{5}}.
\]
\textbf{The error.} The $17$ ``drawings'' are not equally likely. The drawing that starts with $5$ has probability $\dfrac{1}{5}$, while each drawing that starts with one of the other four chips and continues with a specific second chip has probability $\dfrac{1}{5}\cdot\dfrac{1}{4}=\dfrac{1}{20}$. Counting the outcome ``first chip $5$'' as one case among $17$ gives it weight $\dfrac{1}{17}$ instead of $\dfrac{1}{5}$. Weighting correctly gives $4\cdot\dfrac{1}{20}=\dfrac{1}{5}$ again. (Drawing a second chip after a $5$, which never happens in the game, does not change any probability, which is why the simple count over $20$ ordered pairs is valid.)
